\documentclass[11pt]{article}
\usepackage[T1]{fontenc}
\usepackage{lmodern}
\usepackage[margin=1in]{geometry}
\usepackage{amsmath,amssymb,amsthm,mathtools}
\usepackage{microtype}
\usepackage{hyperref}
\hypersetup{colorlinks=true,linkcolor=blue,citecolor=blue,urlcolor=blue,
  pdftitle={From arithmetic dilogarithms to SICs in every dimension},
  pdfauthor={Kaiwen Jin}}
\newtheorem{theorem}{Theorem}
\newtheorem{proposition}[theorem]{Proposition}
\newtheorem{lemma}[theorem]{Lemma}
\newcommand{\C}{\mathbb C}
\newcommand{\Q}{\mathbb Q}
\newcommand{\Z}{\mathbb Z}
\newcommand{\Qbar}{\overline{\mathbb Q}}
\newcommand{\tr}{\operatorname{tr}}
\newcommand{\e}{\mathrm e}
\newcommand{\SF}{\mathcal S}
\title{From arithmetic dilogarithms to SICs in every dimension}
\author{Kaiwen Jin\thanks{Contact: \href{mailto:yui.ui.siki@gmail.com}{yui.ui.siki@gmail.com}.}}
\date{12 October 2026}
\begin{document}
\maketitle
\begin{abstract}
Recent arithmetic results give a construction of symmetric informationally
complete configurations in every finite complex dimension. This note explains
the implication with explicit coordinates and phases. Twisted convolution
produces a trace-one algebraic idempotent. The common-conjugation theorem of
\emph{Stark units for real quadratic fields and reciprocity laws} makes its
nonzero displacement coefficients simultaneously have modulus one, so an
algebraic conjugate is a Hermitian rank-one projector. Its displacement orbit
consists of $d^2$ unit vectors with every distinct squared inner product equal
to $1/(d+1)$. We include the even-dimensional representative convention and
explicit vectors in dimensions two and three. The all-dimension consequence
is already stated in the cited reciprocity preprint; the purpose here is to
spell out its deduction from the arithmetic results.
\end{abstract}

\section{The existence statement and its arithmetic sources}
\label{sec:introduction}

A symmetric informationally complete configuration, or SIC, in $\C^d$ is a
family of $d^2$ unit vectors whose distinct squared inner products are all
$1/(d+1)$. The lines they span are therefore equiangular. The existence of such
a family for every integer $d\geq2$ is the question recorded as AIM~030
\cite{AIM030}. The question imposes no group symmetry on the vectors.

Appleby, Flammia and Kopp \cite{AFK25} gave an arithmetic construction using
special values of the Shintani--Faddeev modular cocycle. Their construction
required both an algebraic conjugation predicted by Stark's conjectures and a
twisted-convolution identity ensuring idempotency. Radchenko and Wheeler
\cite{RW26} proved algebraicity of the corresponding finite quantum
dilogarithm values and a finite pentagon relation. Appleby, Flammia and Kopp
\cite{AFK26} then proved the required convolution identity and supplied an
explicit conversion between the two special-function conventions.
These results produce algebraic rank-one idempotents, called ghost SIC
projectors, but idempotency by itself does not make a matrix Hermitian.

The October preprint \emph{Stark units for real quadratic fields and
reciprocity laws} by Radchenko and Wheeler \cite{REC26} proves the simultaneous conjugation needed for
that last step. Its Proposition~7 states that every algebraic automorphism
switching the two embeddings of the real quadratic field sends all the
nonzero finite dilogarithm values to numbers of modulus one. The same source
already states the SIC consequence in Section~9.2. We give a detailed
deduction in the conventions of \cite{AFK25,AFK26,RW26}, with the coordinate
conversion and all dimensions explicit.

\begin{theorem}[Consequence of the cited arithmetic results]
\label{thm:main}
For every integer $d\geq2$, there exist $d^2$ vectors
$\psi_1,\ldots,\psi_{d^2}$ in $\C^d$ such that
\begin{equation}
\label{eq:target}
 \psi_j^*\psi_j=1,
 \qquad |\psi_j^*\psi_k|^2=\frac1{d+1}\quad(j\ne k).
\end{equation}
\end{theorem}

Here star denotes conjugate transpose, so $v^*w$ is conjugate linear in $v$.
The construction for $d\geq4$ uses a displacement orbit. This supplies
witnesses for the unrestricted assertion~\eqref{eq:target}.
The argument concerns existence; it neither classifies all SICs nor provides
a practical procedure for selecting the algebraic automorphism. Its
arithmetic premises are the exact preprint versions cited below. No
peer-review status is assumed.

The reason modulus one is sufficient is that the displacement coefficients
already come in reciprocal pairs. After one common algebraic conjugation,
the reciprocal of each coefficient becomes its complex conjugate. The
idempotent is then Hermitian, and trace one makes its range a line.
We first fix the displacement and special-function conventions, then derive
the idempotent, and finally prove~\eqref{eq:target}, including $d=2,3$.
All algebraic numbers lie in a fixed algebraic closure $\Qbar\subset\C$.

\section{Displacements and finite arithmetic values}
\label{sec:definitions}

For $d\geq2$, set
\[
 \omega_d=\exp(2\pi i/d),\qquad
 \xi_d=-\exp(\pi i/d),\qquad
 \bar d=\begin{cases}d,&d\text{ odd},\\2d,&d\text{ even}.\end{cases}
\]
In the standard basis $e_0,\ldots,e_{d-1}$ of $\C^d$, define
$Xe_j=e_{j+1\bmod d}$ and $Ze_j=\omega_d^j e_j$. For $p\in\Z^2$, put
\begin{equation}
\label{eq:displacement}
 D_p=\xi_d^{p_1p_2}X^{p_1}Z^{p_2},\qquad
 [p,q]=p_2q_1-p_1q_2.
\end{equation}
The identities $ZX=\omega_dXZ$ and $\xi_d^2=\omega_d$ give
\begin{equation}
\label{eq:weyl}
 D_pD_q=\xi_d^{[p,q]}D_{p+q},\qquad D_p^*=D_{-p}.
\end{equation}
For any fixed representatives of $(\Z/d\Z)^2$, the $d^2$ matrices $D_p$
form a trace-orthogonal basis:
\begin{equation}
\label{eq:trace}
 \tr(D_p)=0\quad(p\notin d\Z^2),\qquad
 \tr(D_p^*D_q)=d\,\delta_{p,q}.
\end{equation}
A nontrivial shift has zero diagonal, and a nontrivial modulation has trace
$\sum_{j=0}^{d-1}\omega_d^{p_2j}=0$, which proves these assertions.
In even dimensions $D_p$ itself can change sign on replacing $p$ modulo $d$.
Its projective class is well defined. Later sums use coefficients for which
the entire summand is representative independent.

We next state the arithmetic result with the normalization needed here.
Let $K$ be a real quadratic field, $I\subset K$ a rank-two lattice, and
$\varepsilon>1$ a totally positive unit satisfying $\varepsilon I=I$.
The prime denotes the other real embedding. Choose an oriented basis
$I=\Z\beta_1+\Z\beta_2$ with
$\beta_1\beta_2'-\beta_1'\beta_2>0$ and $\beta_2$ totally positive.
Put $t=\beta_1/\beta_2$, and let
$A=\left(\begin{smallmatrix}a&b\\c&e\end{smallmatrix}\right)\in
\mathrm{SL}_2(\Z)$ satisfy
$A(\beta_1,\beta_2)^T=\varepsilon(\beta_1,\beta_2)^T$.
Then $A\cdot t=(at+b)/(ct+e)=t$ and $ct+e=\varepsilon$.

Write $\e(z)=\exp(2\pi iz)$ and
$(z;q)_\infty=\prod_{n\geq0}(1-zq^n)$ for $|q|<1$.
The modular quantum dilogarithm is initially the quotient
\begin{equation}
\label{eq:phi}
 \Phi_{A,m,n}(z;t)=
 \frac{(\e(mt+z);\e(t))_\infty}
 {(\e(nA\cdot t+z/(ct+e));\e(A\cdot t))_\infty}
 \qquad(\operatorname{Im}t>0).
\end{equation}
At the real fixed point it means the meromorphic continuation described in
\cite[Section~2.2]{RW26}, rather than direct substitution into either infinite
product. Let $\mu_A$ be the Dedekind eta multiplier, defined in the upper
half-plane by
\[
 \eta(A\cdot t)=\mu_A\sqrt{ct+e}\,\eta(t),\qquad
 \eta(t)=\e(t/24)\prod_{n\geq1}(1-\e(nt)),
\]
with the principal square root. This multiplier is a root of unity.

For a nonzero class in the finite group
$\mathcal G_{I,\varepsilon}=(\varepsilon-1)^{-1}I/I$, choose a representative
$x$ and integers $m,n$ such that
$(\varepsilon^{-1}-1)x/\beta_2=mt+n$. Define
\begin{equation}
\label{eq:E}
 E_{I,\varepsilon}(x)=\mu_A\Phi_{A,m,0}(x/\beta_2;t),
 \qquad E_{I,\varepsilon}(0)=\sqrt\varepsilon>0.
\end{equation}
The algebraicity, reflection and continuation results of
\cite[Theorems~1--2]{RW26}, expressed in these lattice coordinates in
\cite[Section~3 and Appendix~B]{REC26}, make the values nonzero algebraic
numbers and independent of the permitted choices.

\begin{proposition}[Common conjugation, {\cite[Proposition~7]{REC26}}]
\label{prop:common}
For every $K,I,\varepsilon$ as above and every automorphism $\tau$ of $\Qbar$
restricting to the nontrivial automorphism of $K$, one has
\begin{equation}
\label{eq:common}
 |\tau(E_{I,\varepsilon}(x))|=1
 \quad\text{for all }x\in\mathcal G_{I,\varepsilon}\setminus\{0\}.
\end{equation}
\end{proposition}

The quantifier over $\tau$ is shared by all $x$ in this statement. There is
no restriction to a maximal-order lattice or to a fundamental unit. This
allows the lattice and the third power of a unit used below.

\section{A trace-one arithmetic idempotent}
\label{sec:ghost}

Fix $d\geq4$, and set
\begin{equation}
\label{eq:arithmetic}
 \Delta=(d+1)(d-3),\quad
 \rho=\frac{d-1+\sqrt\Delta}{2},\quad
 K=\Q(\sqrt\Delta),\quad I=\Z\rho+\Z,\quad
 \varepsilon=\rho^3.
\end{equation}
We have $(d-2)^2<\Delta<(d-1)^2$, since the respective differences are
$2d-7>0$ and $4$. Thus $\Delta$ is positive and nonsquare. Moreover,
$\rho^2-(d-1)\rho+1=0$ and $\rho'=\rho^{-1}>0$. Hence $\rho$ is a totally
positive integral unit preserving $I$, and $(\rho,1)$ is oriented.
Multiplication by $\rho$ and by $\varepsilon$ is represented respectively by
\begin{equation}
\label{eq:A}
 L=\begin{pmatrix}d-1&-1\\1&0\end{pmatrix},\qquad
 A=L^3=
 \begin{pmatrix}d^3-3d^2+d+1&-d(d-2)\\d(d-2)&1-d\end{pmatrix}.
\end{equation}
In particular,
\begin{equation}
\label{eq:Amod}
 \frac{A-I_2}{d}=
 \begin{pmatrix}d^2-3d+1&2-d\\d-2&-1\end{pmatrix},
 \qquad N=\tr(A)-2=d^2(d-3)>0.
\end{equation}
The fixed point of $A$ is $\rho$ and its lower-left entry is positive.

Here is the precise rank-one instance of the dimension correspondence in
\cite[Definitions~1.28--1.30 and Theorem~4.50]{AFK25}.
The integral form
\[
 Q_d(x,y)=x^2-(d-1)xy+y^2
\]
has discriminant $\Delta$. The tuple $(d,1,Q_d)$ is admissible in
\cite[Definition~1.1]{AFK26}: $1<(d-1)/2$,
$(d^2-1)/(d-1)=d+1$ is an integer, and $\Delta/\operatorname{disc}(Q_d)=1$.
In the notation of \cite{AFK25}, $\rho$ is the $j$th power of the fundamental
totally positive unit, $d=d_{j,1}$, and the conductor of the form is $f=f_j$.
Indeed, if $\Delta_0$ is the field discriminant, this conductor is
$(\rho-\rho')/\sqrt{\Delta_0}$. The source matrix $L_{z,t}$ is exactly $L$,
and its rank-one stabilizer $L_{z,t}^3$ is $A$.

For $p\in\Z^2$ put
\begin{equation}
\label{eq:coordinates}
 S=\begin{pmatrix}0&-1\\1&0\end{pmatrix},\qquad
 u=S^{-1}(A-I_2)p/d,\qquad x_p=\frac{\rho p_2-p_1}{d}\pmod I.
\end{equation}
The vector $u$ is integral by~\eqref{eq:Amod}, with entries
\[
 u_1=(d-2)p_1-p_2,\qquad
 u_2=-(d^2-3d+1)p_1+(d-2)p_2.
\]
Using $A+A^{-1}=(N+2)I_2$ and $\rho^2=(d-1)\rho-1$ gives
\begin{equation}
\label{eq:dictionary-coordinate}
 \frac{(I_2-A^{-1})Su}{N}=\frac pd,\qquad
 \frac{u_1\rho+u_2}{\varepsilon^{-1}-1}=\frac{\rho p_2-p_1}{d}.
\end{equation}
Since $A\equiv I_2\pmod d$, the class $x_p$ belongs to
$\mathcal G_{I,\varepsilon}$; it is zero exactly when $p\in d\Z^2$.
The nonzero branch of the special-function dictionary applies at precisely
these indices. To check its exclusion, put $M=(I_2-A^{-1})S/N$.
Then $Mu=p/d$ and
\[
 M(A^T-I_2)\Z^2=-A^{-1}S\Z^2=\Z^2,
\]
so $u\in(A^T-I_2)\Z^2$ if and only if $p\in d\Z^2$.

For $p\notin d\Z^2$, define
\begin{equation}
\label{eq:nu}
 s_d(p)=d+(1+d)(1+p_1)(1+p_2),\qquad
 \nu_p=\frac{(-1)^{s_d(p)}\xi_d^{-Q_d(p)}}{E_{I,\varepsilon}(x_p)}.
\end{equation}
These are the normalized ghost coefficients of \cite{AFK25}. We spell out
the normalization because an uncancelled multiplier would change the
convolution. Write $\SF_A^r(\rho)$ for the Shintani--Faddeev modular cocycle
at rational column index $r$, in the convention of \cite{AFK25}. On the
nonzero branch, \cite[Proposition~2.1]{AFK26} gives
\begin{equation}
\label{eq:dictionary}
 F_A(u^T)=\mu_A\,\SF_A^{Mu}(\rho)^{-1}
          =E_{I,\varepsilon}(x_p).
\end{equation}
The source phase is
$\phi_p=(-1)^{s_d(p)}\mu_A^{-1}\xi_d^{-Q_d(p)}$, because $f=f_j$ and the
rank is one. Thus $\phi_p\SF_A^{p/d}(\rho)=\nu_p$.
The zero coefficient is separately defined by
\begin{equation}
\label{eq:zero}
 \nu_0=-\varepsilon^{-1/2}
 \qquad\text{{\cite[Equation~(5.65)]{AFK25}}}.
\end{equation}
It is not obtained by inserting the zero class into~\eqref{eq:nu}.
The unconditional reciprocal identity is
\begin{equation}
\label{eq:reciprocal}
 \nu_p\nu_{-p}=1\quad(p\notin d\Z^2)
 \qquad\text{{\cite[Theorem~5.8]{AFK25}}}.
\end{equation}

\begin{lemma}
\label{lem:ghost}
Let $G=\operatorname{diag}(-1,1)$. Then
\begin{equation}
\label{eq:B}
 B=\frac{I_d}{d}+\frac1{d\sqrt{d+1}}
     \sum_{p\ne0\bmod d}\nu_{Gp}D_p
\end{equation}
is an algebraic matrix satisfying $B^2=B$ and $\tr(B)=1$.
The entire summands in this expression are independent of the chosen
integer representatives.
\end{lemma}

\begin{proof}
The proved convolution theorem \cite[Theorem~1.2]{AFK26} permits a shift
$\lambda$ with $\lambda r_{j,m}\equiv-r_{j,m-1}\pmod d$, as specified in
its proof on pages~7--8. At $m=1$ we have $r_{j,0}=0$ and $r_{j,1}=1$,
so $\lambda=0$ is available. With the inverse cocycle identity
\cite[Lemma~2.13]{AFK25}, this gives
\begin{equation}
\label{eq:TCI}
 \sum_q\omega_d^{[p,Lq]}
 \frac{\SF_A^{q/d}(\rho)}{\SF_A^{(q-p)/d}(\rho)}=0
 \qquad(p\ne0\bmod d).
\end{equation}
The representatives in this sum include $0$ and $p$.
We use only this proved shift; no assertion about the entire set of
conjectural shifts is needed.

The twist condition of \cite[Lemma~1.42]{AFK25} is
$\det(G)(2\lambda+d-1+d)\equiv1\pmod{\bar d}$.
For our choices its left side is $1-2d$, so it holds in both parities.
The same lemma gives representative independence of $\nu_{Gp}D_p$.

For the phases $\phi_q$ above, direct calculation yields
\begin{equation}
\label{eq:phase}
 \xi_d^{-[p,q]}\frac{\phi_q}{\phi_{q-p}}
   =(-1)^{1+s_d(p)}\xi_d^{Q_d(p)}\omega_d^{[p,Lq]}.
\end{equation}
Indeed,
\begin{align*}
 s_d(q)+s_d(q-p)&\equiv1+s_d(p)+(1+d)[p,q]\pmod2,\\
 Q_d(q-p)-Q_d(q)&=Q_d(p)-(d-1)[p,q]+2[p,Lq].
\end{align*}
Together with $\xi_d^d=(-1)^{d+1}$ these identities prove~\eqref{eq:phase}.
Multiplying~\eqref{eq:TCI} by the common nonzero phase therefore gives
\begin{equation}
\label{eq:normalizedTCI}
 \sum_q\xi_d^{-[p,q]}\frac{\nu_q}{\nu_{q-p}}=0.
\end{equation}
The zero values in this normalization are exactly~\eqref{eq:zero}.
By~\eqref{eq:reciprocal}, the $q=0,p$ terms are
$\nu_p(\nu_0+\nu_0^{-1})$. Since
$\sqrt\rho+1/\sqrt\rho=\sqrt{d+1}$, we obtain
\[
 \nu_0+\nu_0^{-1}
 =-(\rho^{3/2}+\rho^{-3/2})=-(d-2)\sqrt{d+1}.
\]
Consequently
\begin{equation}
\label{eq:offzero}
 \sum_{q\ne0,p}\xi_d^{-[p,q]}\nu_q\nu_{p-q}
       =(d-2)\sqrt{d+1}\,\nu_p.
\end{equation}

Put $F=\sum_{p\ne0}\nu_pD_{G^{-1}p}$. Since $\det(G^{-1})=-1$,
\eqref{eq:weyl} gives the phase $\xi_d^{-[p,q]}$ in the coefficient of
$D_{G^{-1}p}$ in $F^2$. Formula~\eqref{eq:offzero} computes that
coefficient. The identity coefficient is $d^2-1$, by~\eqref{eq:reciprocal}.
Thus
\begin{equation}
\label{eq:F}
 F^2=(d^2-1)I_d+(d-2)\sqrt{d+1}\,F.
\end{equation}
Representative independence justifies the changes of index also when $d$ is
even. Substitution in $B=I_d/d+F/(d\sqrt{d+1})$ proves $B^2=B$.
Equation~\eqref{eq:trace} gives $\tr(B)=1$.
All its entries are algebraic, by the arithmetic algebraicity theorem and
\eqref{eq:nu}. No Hermiticity assertion has been used.
\end{proof}

\section{Hermitian conjugation and the vector family}
\label{sec:hermitian}

The nontrivial automorphism of $K$ extends to an automorphism $\tau$ of
$\Qbar$. For completeness, order partial field embeddings into $\Qbar$
extending that automorphism by extension of their domains. Unions of chains
are embeddings, so Zorn's lemma gives a maximal one. If its domain omits an
algebraic number, a root of the transformed minimal polynomial extends the
embedding, a contradiction. Its domain is therefore $\Qbar$. Its image is
algebraically closed: pull back a polynomial with coefficients in the image,
take a root, and map that root. An algebraically closed subfield of $\Qbar$
containing $\Q$ contains every root of every rational polynomial, so the
image is all of $\Qbar$. This proves the required automorphism extension.
It may be selected separately for each $d$.

Apply Proposition~\ref{prop:common} to this one $\tau$ and all $x_{Gp}$.
Roots of unity map to roots of unity, so
\begin{equation}
\label{eq:c}
 c_p=\tau(\nu_{Gp}),\qquad |c_p|=1,\qquad
 c_pc_{-p}=1,\qquad c_{-p}=\overline{c_p}.
\end{equation}
We must also conjugate the displacement matrices and the normalization.
There are $k$ coprime to $\bar d$ and $h\in\{1,-1\}$ with
\[
 \tau(\xi_d)=\xi_d^k,\qquad
 \tau(\sqrt{d+1})=h\sqrt{d+1}.
\]
Since $X$ has rational entries and $\tau(Z)=Z^k$, entrywise conjugation gives
\begin{equation}
\label{eq:P}
 U_p=\tau(D_p)=D_{(p_1,kp_2)},\qquad
 P=\tau(B)=\frac{I_d}{d}+\frac{h}{d\sqrt{d+1}}
           \sum_{p\ne0}c_pU_p.
\end{equation}
In particular $U_{-p}=U_p^*$, and $p\mapsto(p_1,kp_2)$ permutes the
classes modulo $d$. The full summands $c_pU_p$ remain representative
independent. Reindexing their adjoints by $p\mapsto-p$ and
using~\eqref{eq:c} proves $P^*=P$. This argument includes self-inverse
classes in even dimensions. It does not require $\tau$ to commute with
complex conjugation on arbitrary algebraic numbers.

Applying a field automorphism to the algebraic equations in
Lemma~\ref{lem:ghost} gives $P^2=P$ and $\tr(P)=1$.
Every idempotent has the direct sum decomposition
$\C^d=\operatorname{ran}P\oplus\ker P$, and is diagonal with ones on its
range and zeros on its kernel in an adapted basis. Therefore
$\operatorname{rank}P=\tr(P)=1$. Hermiticity makes the kernel orthogonal to
the range. Choose a unit vector $\psi$ spanning that range. Then
\begin{equation}
\label{eq:projector}
 P=\psi\psi^*,\qquad v^*Pv=\|Pv\|^2\geq0\quad(v\in\C^d).
\end{equation}

For every nonzero class $q$ modulo $d$, trace orthogonality and the
permutation in~\eqref{eq:P} imply
\begin{equation}
\label{eq:overlap}
 |\tr(PD_q^*)|=\frac1{\sqrt{d+1}},\qquad
 |\psi^*D_q\psi|^2=\frac1{d+1}.
\end{equation}
Changing an integer representative only multiplies the relevant
displacement by a phase of modulus one. The possible sign $h$ likewise
disappears from the absolute value.
For one representative of each $q\in(\Z/d\Z)^2$ define $\psi_q=D_q\psi$.
All these vectors have unit norm. If $q\ne r$, then $D_q^*D_r$ is a unit
phase times $D_{r-q}$, and hence
\[
 |\psi_q^*\psi_r|^2=|\psi^*D_{r-q}\psi|^2=\frac1{d+1}.
\]
Since this number is less than one, no two vectors are parallel. We have
therefore obtained exactly $d^2$ vectors satisfying~\eqref{eq:target}
for every $d\geq4$.

The two remaining dimensions have elementary witnesses. For $d=2$, set
$\zeta=\exp(2\pi i/3)$ and take
\[
 v_0=(1,0),\qquad
 v_{j+1}=\bigl(1/\sqrt3,\sqrt{2/3}\,\zeta^j\bigr),\quad j=0,1,2.
\]
Their norms are one. The squared overlaps with $v_0$ equal $1/3$, and for
$j\ne k$ the other squared overlaps equal
$|1+2\zeta^{k-j}|^2/9=1/3$, since
$\operatorname{Re}\zeta=\operatorname{Re}\zeta^2=-1/2$.

For $d=3$, take $v=(0,1,-1)/\sqrt2$ and the nine vectors
$X^pZ^qv$, with $p,q=0,1,2$ and the dimension-three generators defined
in Section~\ref{sec:definitions}. For $q=1,2$,
\[
 v^*Z^qv=(\omega_3^q+\omega_3^{2q})/2=-1/2.
\]
For $p=1$ the overlap $v^*X^pZ^qv$ is $-\omega_3^q/2$, and for $p=2$
it is $-\omega_3^{2q}/2$. Thus every nontrivial shift or modulation has
overlap of modulus $1/2$. Weyl multiplication reduces every distinct pair
in the orbit to one of these overlaps, so its squared inner product is
$1/4$. Together with the $d\geq4$ construction, these two cases prove
Theorem~\ref{thm:main} for every integer $d\geq2$.

\paragraph{AI assistance.}
OpenAI Codex assisted with source retrieval, mathematical derivations,
checking, and preparation of this note. The cited reciprocity preprint also
discloses substantial LLM assistance in its own preparation. These disclosures
do not replace mathematical verification or independent review.

\begin{thebibliography}{9}
\bibitem{AFK25}
M. Appleby, S. T. Flammia and G. S. Kopp,
\emph{A Constructive Approach to Zauner's Conjecture via the Stark Conjectures},
arXiv:2501.03970v2 (17 March 2025).
\url{https://arxiv.org/abs/2501.03970v2}.

\bibitem{RW26}
D. Radchenko and C. Wheeler,
\emph{Real quadratic fields and finite quantum dilogarithms I},
arXiv:2609.21892v2 (29 September 2026).
\url{https://arxiv.org/abs/2609.21892v2}.

\bibitem{AFK26}
M. Appleby, S. T. Flammia and G. S. Kopp,
\emph{The twisted convolution identity and ghost $r$-SICs from finite
quantum dilogarithms}, arXiv:2609.39192v1 (30 September 2026).
\url{https://arxiv.org/abs/2609.39192v1}.

\bibitem{REC26}
Danylo Radchenko and Campbell Wheeler,
\emph{Stark units for real quadratic fields and reciprocity laws},
Hexagon 2610.00186v1 (10 October 2026), especially Proposition~7 and
Section~9.2.
\url{https://hexagonmath.org/2610.00186v1}.

\bibitem{AIM030}
AIM problem catalogue (\texttt{MColbrook/AIM}),
\emph{030 --- SIC measurements in every finite dimension},
statement at revision \texttt{c9929805cece044705abaaebe1cebddac1ba11a6}.
\href{https://github.com/MColbrook/AIM/blob/c9929805cece044705abaaebe1cebddac1ba11a6/problems/030-sic-povms.md}{Frozen AIM 030 statement (GitHub)}.
\end{thebibliography}
\end{document}
